\section{Support and Zero-Localization}
\label{sec:zero-localization}

The most deceptive obstacle in the formalization is support.  Mathematicians
habitually speak about a chart representative only where the chart is defined.
Lean does not.  A chart representative is a total function, and its support is
a global statement about that total function.  The central discovery of this
part of the development is that the problem cannot be solved by asking for a
stronger support hypothesis on the same representative.  The representative
itself has to be changed.

\subsection{The false support theorem}

Suppose \(\rho\) is a partition coefficient and \(\omega\) is a form.  A paper
proof may say that the coordinate representative of \(\rho\omega\) is
supported in a coordinate box \(Q\).  In Lean, the corresponding ordinary
representative is
\[
\code{ManifoldForm.transitionPullbackInChart}.
\]
It has a value at every coordinate point.  Outside the chart source, target, or
overlap, those values are artifacts of totalization.  Thus a theorem of the
form
\[
\tsupp(\code{transitionPullbackInChart}(\rho\omega)) \subseteq Q
\]
may be too strong or simply unrelated to the informal argument.  It asks Lean
to control the chart representative even at points where the coordinate
description should not be considered part of the local proof.

This problem explains several early certificate fields in the development.  It
is easy to request a support assumption from the user.  It is harder, and more
mathematically faithful, to prove the right theorem from compact support and
partition support control.  The right theorem is not about the unrestricted
total representative.

One can view the issue as a small counterexample pattern.  Suppose a chart
expression is mathematically intended only on a set \(U\), but Lean represents
it by some total function \(f : E \to V\).  Even if \(f\) is perfectly
controlled on \(U\), its values on \(E\setminus U\) may enlarge
\(\operatorname{support}(f)\).  A theorem about \(\operatorname{support}(f)\)
is therefore not the same as a theorem about the local representative on
\(U\).  The zero-localized representative replaces \(f\) by
\[
  y\mapsto
  \begin{cases}
    f(y), & y\in U,\\
    0, & y\notin U,
  \end{cases}
\]
so that global support again expresses the intended local support.  This is a
change of object, not a notational convenience.

\begin{remark}[The false support theorem]
The tempting theorem is:
\[
  \tsupp(\code{transitionPullbackInChart}(\rho\omega))\subseteq Q.
\]
The formalization treats this as the wrong target.  The values of
\code{transitionPullbackInChart} outside the transition source are present
because Lean functions are total, not because the local coordinate proof has
mathematical content there.  Strengthening the hypotheses to control those
values would make the final theorem depend on an artifact of representation.
The corrected target is the same inclusion for the zero-localized
representative, together with an agreement theorem on the transition source.
\end{remark}

\subsection{Zero outside the chart domain}

The project introduces a zero-localized representative:
\[
\code{ManifoldForm.transitionPullbackInChartZero}.
\]
It agrees with \code{transitionPullbackInChart} on the relevant chart source
and is zero outside.  The basic API proves:
\begin{itemize}
  \item support is contained in the source;
  \item support is contained in the target;
  \item support is contained in the overlap;
  \item the zero representative agrees with the ordinary representative on
  the chart source.
\end{itemize}
The core names are:
\begin{lstlisting}[language=Lean]
transitionPullbackInChartZero_support_subset_source
transitionPullbackInChartZero_support_subset_target
transitionPullbackInChartZero_support_subset_overlap
transitionPullbackInChartZero_eqOn_transitionPullbackInChart
\end{lstlisting}

This changes the meaning of support control.  Instead of proving that a
totalized chart function behaves well outside its mathematical domain, the
proof makes the representative zero there by definition and proves support
properties for the resulting function.  The point is not simply that zero
extension is available; it is that the zero extension is the formal object
corresponding to the informal phrase ``the local chart representative''.

The equality theorem is as important as the support theorem.  The zero
representative must agree with the ordinary representative on the region where
local integration is performed; otherwise the theorem would prove Stokes for a
different object.  This is why the API includes both support inclusions and
agreement theorems such as
\code{transitionPullbackInChartZero_eqOn_transitionPullbackInChart}.

The zero-localization API therefore has a three-part mathematical shape:
\[
\begin{array}{ll}
\text{definition} & \text{make the off-source representative zero},\\
\text{support theorem} & \text{prove source, target, and overlap support},\\
\text{agreement theorem} & \text{recover the ordinary chart expression on the
source}.
\end{array}
\]
All three are needed.  The definition alone would give a different
representative; the agreement theorem makes it harmless for integration on
the selected box; the support theorem makes artificial-face cancellation
available.

\subsection{Localized forms}

The global proof uses localized forms \(\rho\omega\).  Zero-localization has to
interact with localization.  The theorem
\[
\code{ManifoldForm.transitionPullbackInChartZero_localizedForm}
\]
rewrites the zero-localized representative of a localized form as the product
of the coefficient representative and the zero-localized representative of the
base form.  From this, the project proves support inclusions such as
\[
\code{transitionPullbackInChartZero_localizedForm_tsupport_subset_coefficient}
\]
and
\[
\code{transitionPullbackInChartZero_localizedForm_tsupport_subset_form}.
\]

The decisive bridge is:
\begin{lstlisting}[language=Lean,basicstyle=\ttfamily\scriptsize]
transitionPullbackInChartZero_localizedForm_tsupport_subset_halfSpaceSupportBox_of_support_subset_K
\end{lstlisting}
It formalizes the following argument:
\[
\begin{array}{c}
\supp(\omega)\subseteq K,\\
\tsupp(\rho)\cap K\text{ is controlled by the selected chart box},\\
\text{the coordinate image of the active carrier is controlled},\\
\hline
\tsupp(\text{zero-localized representative of }\rho\omega)
  \subseteq \text{selected half-space support box}.
\end{array}
\]
This is the support theorem the half-space local Stokes layer needs.

\begin{lemma}[Zero-localized support bridge]
\label{lem:zero-support}
Let \(K\) be a compact carrier containing \(\supp(\omega)\).  Let \(\rho\) be a
localization coefficient whose topological support on \(K\) is contained in a
coordinate carrier, and suppose that the coordinate image of that carrier is
contained in a half-space support box \(Q(a,b)\).  Then the zero-localized chart
representative of \(\rho\omega\) has topological support contained in
\(Q(a,b)\).
\end{lemma}

This lemma is the formal replacement for the informal statement that
``\(\rho\omega\) is supported in the selected box.''  The latter statement is
ambiguous in a total-function setting; the former is a precise theorem about
the representative actually used by the local Stokes theorem.

At the level of Lean proof dependencies, the bridge decomposes as follows:
\[
\begin{array}{c}
\tsupp(\operatorname{ZeroRep}(\rho\omega))
  \subseteq
  \tsupp(\rho\text{ in chart}),\\[2mm]
\tsupp(\operatorname{ZeroRep}(\rho\omega))
  \subseteq
  \tsupp(\operatorname{ZeroRep}(\omega)),\\[2mm]
\tsupp(\operatorname{ZeroRep}(\omega))
  \subseteq \operatorname{chartImage}(K),\\[2mm]
\tsupp(\rho\text{ in chart})\cap \operatorname{chartImage}(K)
  \subseteq Q(a,b),\\[2mm]
\hline
\tsupp(\operatorname{ZeroRep}(\rho\omega))\subseteq Q(a,b).
\end{array}
\]
The first two inclusions are product-support facts.  The third is the global
compact-support input transported through the zero representative.  The last
line is the selected-refinement support control.  Their combination is exactly
the support bridge displayed above.

There are two support operations in the proof.  Ordinary support is algebraic:
it records where a function is nonzero.  Topological support is the closure of
ordinary support and is the right object for face-disjointness and vanishing
integrals.  The development therefore does not stop at ordinary nonzero-set
control.  The final local Stokes input is a \(\tsupp\)-level inclusion into a
selected half-space support region.  That region is designed to allow contact
with the true boundary face \(x_0=0\) while strictly avoiding the artificial
faces of the coordinate box.

\subsection{Selected smooth refinement}

The final route applies the preceding bridge to the canonical smooth
refinement selected by the first-principles theorem.  The theorem is:
\begin{lstlisting}[language=Lean]
CoverIndexedZeroCompactRepresentedStokesIntrinsicInput.
  zeroSupportOfSelectedSmoothRefinement
\end{lstlisting}
states that for every selected boundary index \(i\) and active refined
half-space piece \(q\), the topological support of the zero-localized
representative of the refined localized form is contained in the selected
half-space support box:
\[
\tsupp
  \left(
    \operatorname{ZeroRep}
    \bigl(\rho_{i,q}\omega\bigr)
  \right)
  \subseteq
  \code{halfSpaceSupportBox}(a_{i,q},b_{i,q}).
\]
Here \(\rho_{i,q}\) is not a user-supplied coefficient.  It is the coefficient
generated by the selected smooth refinement.

This point is crucial for the first-principles theorem.  If \(\rho_{i,q}\)
were arbitrary, the zero-support theorem would be false.  It is true because
\(\rho_{i,q}\) is built from the selected boundary partition and the selected
smooth subpartition, both of which carry support-control data from the
compactness construction.

The proof uses a two-stage localization identity.  The refined coefficient is
the product of the selected boundary partition coefficient and the smooth
subpartition coefficient.  The proof rewrites
\[
  \rho_{i,q}\omega
  =
  \theta_{i,q}(\rho_i\omega)
\]
and then uses the support of \(\rho_i\omega\) to place the form in the
boundary active carrier, while the support of \(\theta_{i,q}\) places the
coordinate representative inside the refined half-space box.

\begin{proof}[Proof idea for the selected theorem]
The proof first shows that the boundary-localized form \(\rho_i\omega\) is
supported in the boundary active carrier.  This uses the algebraic fact that
support of a localized form is contained in the support of the coefficient and
in the support of the original form, together with
\(\supp(\omega)\subseteq K\).  Next, the smooth subpartition coefficient
\(\theta_{i,q}\) is known, on that active carrier, to be supported inside the
selected refined half-space box.  The zero-localized product formula transfers
these two support facts to the chart representative.  The proof is carried out
at the level of topological support for the coefficient and the zero base
representative, using the product-support inclusions
\[
  \tsupp(fg)\subseteq\tsupp(f),
  \qquad
  \tsupp(fg)\subseteq\tsupp(g).
\]
Thus the theorem delivered to the half-space layer is already the required
\(\tsupp\)-inclusion; it is not a fragile ordinary-support statement later
patched by hand.
\end{proof}

\subsection{Ordinary support versus topological support}

Another formal distinction is ordinary support versus topological support.
Many algebraic arguments first show that if a function is nonzero at a point,
then that point lies in a desired set.  But artificial-face vanishing is a
closed-support statement.  The theorem must control
\[
  \tsupp(f)=\overline{\{x\mid f(x)\ne0\}}.
\]
In the half-space box proof the target support region is not a naive closed
box.  The set \(\code{halfSpaceSupportBox}(a,b)\) allows the support to meet
the true lower normal face, but uses strict inequalities in the directions
corresponding to artificial faces.  This is exactly the geometry needed by the
face-cancellation theorem: it proves that the artificial face set is disjoint
from the topological support.  Consequently the proof is organized to produce
\(\tsupp\)-statements directly for the localized coefficients and
zero-localized representatives.

Closed-set upgrades still appear in the surrounding support infrastructure in
the familiar form
\[
  \operatorname{support}(f)\subseteq C,\quad C\text{ closed}
  \quad\Longrightarrow\quad
  \tsupp(f)\subseteq C.
\]
But the final half-space input is more specific: it is a topological-support
inclusion into a region chosen so that artificial faces are avoided and the
true boundary face is retained.  This is why the support theorem is not merely
``the coefficient is zero on the face.''  It is a coordinated package of
zero-extension, product support, compact carrier support, and half-space box
geometry.

\subsection{Contribution to the global theorem}

Zero-localized support is the bridge between compact finite data and local
Stokes.  The support-controlled partition and smooth refinement produce
coefficient support information.  Zero-localization turns it into support
information for the coordinate form representative.  The half-space layer then
uses that support information to remove artificial faces.  Without this bridge,
the final theorem would have to expose an unnatural support assumption for
ordinary total chart representatives.  With it, the first-principles theorem
requires no such field.
